\chapter*{\centering \Large Abstract}

\noindent
He et al. \citep{he2016eals} introduce eALS, an element-wise Alternating
Least Squares learner for matrix factorization on implicit feedback that
removes the two main bottlenecks of earlier whole-data ALS solutions: the
cubic dependency on the number of latent factors $K$ inherited from matrix
inversion, and the uniform weighting of unobserved (missing) entries, which
the authors replace with an item-popularity-aware scheme. The paper proves
that the resulting per-factor update rule is embarrassingly parallel, but
ships no distributed implementation to back that claim up.

\medskip
This report closes that gap. eALS is implemented from scratch in PySpark,
following Algorithm~1 of the paper equation by equation and parallelized by
partitioning the user (respectively item) index space across executors,
along the lines the paper's Section~4.2.1 describes but does not build. The
implementation is evaluated on the Amazon ``Movies and TV'' implicit-feedback
dataset, one of the two datasets used in the original paper, filtered with
the same 10-core protocol and evaluated with the same leave-one-out,
Hit-Ratio/NDCG methodology. We reproduce the paper's two main experiments:
the impact of the missing-data weighting parameters $c_0$ and $\alpha$, and
the convergence behaviour of eALS against a uniform-weight ablation. We also
add a dedicated scalability study, something the paper only discusses
qualitatively, measuring training time per iteration against the number of
latent factors, the degree of Spark parallelism, and the size of the
training set. On this dataset eALS scales close to linearly in $K$ (roughly
$\times 11$ from $K{=}8$ to $K{=}128$, against the $\times 16$ a cubic
dependency would need), which confirms the paper's central efficiency
claim. It also exposes a practical limit the paper's single-thread
benchmarks never had to confront: on a single machine, Spark's
task-scheduling overhead outweighs the benefit of more partitions beyond
roughly four to eight. All code, the full experimental log, and a
small-data notebook that exercises the entire pipeline end to end are
included in the appendix.
